\begingroup

\setcounter{section}{0}
\setcounter{theorem}{0}
\setcounter{lemma}{0}
\setcounter{corollary}{0}
\setcounter{proposition}{0}
\setcounter{remark}{0}
\setcounter{equation}{0}
\section{A concrete dimension-independent local
theorem}\label{U:a-concrete-dimension-independent-local-theorem}

Let \(a\in\mathbb R^n\) be a unit vector, \(b\in\mathbb R\), and write
\[
 P=\{x:a^Tx=b\},\qquad S=a^\perp,\qquad P_S=I-aa^T.
\] The normal direction is fixed; the dimension \(n-1\) of the tangent
space is unrestricted.

\phantomsection\label{U:one}\textbf{Theorem U.1 (One exceptional hyperplane).} Let
\(F:\mathbb R^n\to\mathbb R\) be finite and convex, with unique
minimizer \(x_*\in P\). Let \(K=\{F\le F(x_0)\}\) be compact. On an open
convex neighborhood of \(K\), assume:

\begin{enumerate}
\def\labelenumi{\arabic{enumi}.}
\tightlist
\item
  \(F\) is \(\mu\)-strongly convex for some \(\mu>0\).
\item
  All subgradients have the same tangential projection,
  \[P_Sg=A(x)\qquad(g\in\partial F(x)),\] where \(A\) is Lipschitz
  continuous on that neighborhood.
\item
  \(F\) is locally \(C^{1,1}\) away from \(P\).
\end{enumerate}

Then every infinite admissible full-BFGS run starting at \(x_0\)
satisfies \[
 \boxed{F(x_k)\to F_*,\qquad x_k\to x_*.}
\] The hyperplane itself may include differentiable points; it is
allowed to contain failures of gradient-Lipschitz regularity as well as
actual kinks. No uniform bound on the full gradient-Lipschitz constants
as the hyperplane is approached is required.

If these hypotheses hold only in a neighborhood of a unique minimizer,
some \(\varepsilon>0\) gives the same conclusion for every start with
\(F(x_0)<F_*+\varepsilon\). The same basin works for every SPD
initialization and every fixed valid line-search pair. Compact sublevel
trapping supplies this local statement: choose a low sublevel contained
in the structural neighborhood, then use Armijo monotonicity.

The assumptions concern the objective. In particular, the claim that a
part of \(B_k=H_k^{-1}\) has a controlled growth rate will be
\textbf{derived}, not assumed.

\subsection{1.1 A simple family and a weaker-curvature
version}\label{U:a-simple-family-and-a-weaker-curvature-version}

An immediate family is \[
 \boxed{F(u,v)=\psi(u,v)+h(v),\qquad u\in\mathbb R^{n-1},\quad v\in\mathbb R,}
\] where \(\psi\) is convex with locally Lipschitz gradient, \(h\) is
finite convex and locally \(C^{1,1}\) on \(\mathbb R\setminus\{0\}\),
\(F\) is strongly convex on the relevant neighborhood, and the minimizer
has \(v_*=0\). Here \(A=\nabla_u\psi\), irrespective of the singular
normal derivative of \(h\). The objective and the initial matrix may be
fully coupled.

Full strong convexity can be weakened for this composite family. It is
enough that \[
 \psi(u,v)-\frac{\mu}{2}\|u\|^2\quad\hbox{is convex},\qquad \mu>0,
\] that the relevant positive sublevels of \(F\) are strictly convex,
and that the minimizer is unique with \(v_*=0\). A compact initial
sublevel is still required. The proof of this extension is given in
Section U.6. It needs positive curvature only in the smooth tangential
directions, not a positive lower bound in the normal direction.

This is not a claim that all convex nonsmooth functions admit such a
splitting. Nor does a common hyperplane through the optimizer describe
arbitrary curved switching geometry.


\setcounter{section}{1}
\setcounter{theorem}{0}
\setcounter{lemma}{0}
\setcounter{corollary}{0}
\setcounter{proposition}{0}
\setcounter{remark}{0}
\setcounter{equation}{0}
\section{A forced affine-critical cluster-point
theorem}\label{U:a-forced-affine-critical-cluster-point-theorem}

The method also fits the affine-criticality program directly.

\phantomsection\label{U:parallel}\textbf{Theorem U.2 (Parallel singular interfaces).} Let \(F\) be finite
convex with attained minimum and compact sublevel \(K\). Assume \(F\) is
\(\mu\)-strongly convex on an open convex neighborhood of \(K\). Let
\(a\) be a unit vector and \(E\subset\mathbb R\) closed. On that
neighborhood suppose:

\begin{itemize}
\tightlist
\item
  the singular set is \(\Sigma=\{x:a^Tx\in E\}\), every point of
  \(\Sigma\) is genuinely nondifferentiable, and \(F\) is locally
  \(C^{1,1}\) off \(\Sigma\);
\item
  \(P_S\partial F(x)=\{A(x)\}\) for a Lipschitz field \(A\), where
  \(S=a^\perp\).
\end{itemize}

Then the cluster set of every infinite admissible full-BFGS run in \(K\)
contains a point \(p\) with \[
 \boxed{0\in\operatorname{aff}\partial F(p).}
\] There may be several or countably many parallel interfaces, subject
to the closed-set and regularity assumptions above. No finite
affine-critical set is assumed.

At a point of \(\Sigma\), the subdifferential is a nontrivial compact
segment parallel to \(a\), because every subgradient has the same \(S\)
projection. Hence \[
 \operatorname{aff}\partial F(p)=A(p)+\mathbb R a,
 \qquad
 0\in\operatorname{aff}\partial F(p)\iff A(p)=0.
\] At a smooth point, affine criticality is ordinary stationarity.

Consequently, adding \[
 0\in\operatorname{aff}\partial F(p)\implies0\in\partial F(p)
\] on \(K\), or on a sufficiently low trapped neighborhood, gives
convergence to the minimum value and, under uniqueness, convergence of
the whole sequence.

Theorem U.1 is slightly different from a literal specialization of Theorem U.2: it allows the exceptional hyperplane to contain smooth but not
locally \(C^{1,1}\) points. Having \(x_*\in P\) makes every zero of
\(A\) on \(P\) optimal, so genuine nondifferentiability at every
exceptional point is unnecessary there.


\setcounter{section}{4}
\setcounter{theorem}{0}
\setcounter{lemma}{0}
\setcounter{corollary}{0}
\setcounter{proposition}{0}
\setcounter{remark}{0}
\setcounter{equation}{0}
\section{Proof of the parallel-interface and single-interface
results}\label{U:proof-of-the-parallel-interface-and-single-interface-results}

We give a common argument that covers Theorems 1 and 2.

\subsection{5.1 Preliminary bounds and a contradiction
set}\label{U:preliminary-bounds-and-a-contradiction-set}

Strong convexity and convexity give \[
 d_k-(F_k-F_{k+1})\ge\frac\mu2\|s_k\|^2,
 \qquad t_k\ge\mu\|s_k\|^2.\tag{U.10}
\] Thus \(\sum\|s_k\|^2<\infty\) and \(s_k\to0\). Let \(\Gamma\) be the
compact cluster set. Every point of \(\Gamma\) has the common limit
value \(\ell\).

For Theorem U.2, assume \(\Gamma\) has no affine-critical point. Then
\(\ell>F_*\), so it contains no stationary point. For Theorem U.1, assume
directly \(\ell>F_*\).

Write \(r(x)=a^Tx\) in Theorem U.2, and \(r(x)=a^Tx-b\) in Theorem U.1. Let
\(E\) be the exceptional scalar set; in Theorem U.1 use \(E=\{0\}\).

The common tangential projection and its Lipschitz constant \(L_A\)
imply \[
 \|P_S y_k\|^2=\|A(x_{k+1})-A(x_k)\|^2
 \le L_A^2\|s_k\|^2\le\frac{L_A^2}{\mu}t_k.\tag{U.11}
\] Hence Theorem~\ref{M:residualtheorem}'s smooth-subspace estimate holds.

\subsection{5.2 Why tangential descent cannot disappear on the
exceptional
set}\label{U:why-tangential-descent-cannot-disappear-on-the-exceptional-set}

For Theorem U.2, at an exceptional point the nontrivial subgradient
segment has affine hull \(A(p)+\mathbb R a\). Thus \(A(p)=0\) would make
it affine critical, contrary to the assumption on \(\Gamma\).

For Theorem U.1, if \(p\in P\) and \(g\in\partial F(p)\), strong
monotonicity with \(0\in\partial F(x_*)\) gives \[
 A(p)^T(p-x_*)=g^T(p-x_*)\ge\mu\|p-x_*\|^2.
\] Thus \(A(p)=0\) on \(P\) is possible only at \(x_*\). Such a point
cannot belong to a positive limiting level.

In either case the compact set \[
 Z=\{p\in\Gamma:A(p)=0\}
\] has scalar image \(r(Z)\) disjoint from the closed set \(E\). If
\(Z\) is nonempty, these sets have positive separation
\(\delta=\operatorname{dist}(r(Z),E)>0\). If \(E\) is empty there is no
singular region and we may take the cutoff below identically one; if
\(Z\) is empty it may be identically zero.

Choose a scalar Lipschitz function \(\tau:\mathbb R\to[0,1]\) which is
zero where \(\operatorname{dist}(r,E)\le\delta/4\), and one where
\(\operatorname{dist}(r,E)\ge\delta/2\). Such a cutoff is obtained by
composing a piecewise-linear cutoff with the distance function. In
particular \(\tau=1\) on \(r(Z)\). The cutoff is an analytical object
chosen from \(F\) and the hypothetical cluster set, not an algorithmic
input.

\subsection{5.3 A frame on the entire cluster
set}\label{U:a-frame-on-the-entire-cluster-set}

Let the columns of \(U\) be an orthonormal basis of \(S\), and set \[
 \boxed{V(x)=[\,U,\ \tau(r(x))a\,].}\tag{U.12}
\] Its tangential columns are always present. Its normal column is
switched off near singular interfaces and restored away from them.

At differentiable points, \[
 \|V(x)^T\nabla F(x)\|^2
 =\|A(x)\|^2+\tau(r(x))^2|a^T\nabla F(x)|^2.
\] On \(\Gamma\setminus Z\) the first term is positive. On \(Z\), the
cutoff equals one and the full gradient is nonzero. The displayed
expression extends continuously across singular interfaces because the
cutoff vanishes on a whole strip around them; gradients are bounded and
\(A\) is continuous. Compactness therefore provides \(m>0\) and a
buffered neighborhood of all of \(\Gamma\) in which visibility holds.

The entire late trajectory lies in this neighborhood because its
distance to its compact cluster set tends to zero. There is no selection
of one cluster point followed by an assumption that the sequence
converges to it.

\subsection{5.4 The new secant term is
bounded}\label{U:the-new-secant-term-is-bounded}

The tangential part follows from (11). If \(\tau(r(z))=0\), there is no
normal contribution.

If \(\tau(r(z))>0\), the endpoint is a fixed positive distance from the
exceptional scalar set. For sufficiently close endpoints, the entire
segment lies in a slightly smaller smooth strip. On a compact buffered
subset of that strip, a finite convex-chart cover gives a common local
gradient-Lipschitz constant. Convex cocoercivity on each sufficiently
small convex chart yields \[
 \|\nabla F(z)-\nabla F(x)\|^2\le L_\delta\,y^Ts.
\] Therefore \[
 \|V(z)^Ty\|^2\le C_0t\tag{U.13}
\] for every sufficiently late accepted pair. This uses local smoothness
away from the exceptional set; \(L_\delta\) need not stay bounded as
\(\delta\downarrow0\).

\subsection{5.5 The only transport error is tangential and
inexpensive}\label{U:the-only-transport-error-is-tangential-and-inexpensive}

For an endpoint pair put \[
 \beta=\begin{cases}
 [\tau(r(z))-\tau(r(x))]/(a^Ts),&a^Ts\ne0,\\
 0,&a^Ts=0.
 \end{cases}
\] Then \(|\beta|\le\operatorname{Lip}(\tau)\). With \(e_n\) the last
frame-coordinate vector, \[
 \boxed{
 V(z)=V(x)+s\,\beta e_n^T-\beta P_Ss\,e_n^T.
 }\tag{U.14}
\] Thus (6) holds with \(T=I\) and \[
 R=-\beta P_Ss\,e_n^T\in S,
 \qquad
 \|R\|_F^2\le\operatorname{Lip}(\tau)^2\|P_Ss\|^2
 \le\frac{2\operatorname{Lip}(\tau)^2}{\mu}d.\tag{U.15}
\] All conditions of Theorem~\ref{M:residualtheorem} hold. The contradiction proves Theorem U.2,
and proves \(\ell=F_*\) for Theorem U.1. Compactness and uniqueness give
whole-sequence convergence. \(\square\)


\setcounter{section}{5}
\setcounter{theorem}{0}
\setcounter{lemma}{0}
\setcounter{corollary}{0}
\setcounter{proposition}{0}
\setcounter{remark}{0}
\setcounter{equation}{0}
\section{Only tangential strong convexity is needed for a composite
subclass}\label{U:only-tangential-strong-convexity-is-needed-for-a-composite-subclass}

Consider \(F(u,v)=\psi(u,v)+h(v)\) under the weaker hypotheses in
Section U.1.1. Set \(S=\mathbb R^{n-1}\times\{0\}\). Since
\(\psi-\mu\|u\|^2/2\) and \(h\) are convex, \[
 d_k-(F_k-F_{k+1})\ge\frac\mu2\|P_Ss_k\|^2.\tag{U.16}
\] This is the bound needed for the transport residual, even without a
positive normal curvature bound.

At differentiable accepted points, \[
 P_Sy=P_S[\nabla\psi(z)-\nabla\psi(x)].
\] Convex \(L\)-smoothness of \(\psi\) on the compact neighborhood, and
monotonicity of the scalar subgradient of \(h\), give \[
 \|P_Sy\|^2
 \le L[\nabla\psi(z)-\nabla\psi(x)]^Ts
 \le Lt.\tag{U.17}
\] This replaces (11). On the exceptional plane \(v=0\), the restriction
of \(F\) is a differentiable strongly convex function of \(u\), uniquely
minimized at \(u_*\). Thus the needed nonvanishing of \(A=\nabla_u\psi\)
on a nonoptimal cluster set still holds.

Finally, strict convexity of the positive sublevels supplies \(s_k\to0\)
in place of full strong convexity. The proof is dimension independent:
if a subsequence had \(x_k\to x\), \(x_{k+1}\to z\ne x\), \(g_k\to g\),
and \(F_k\to\ell>F_*\), then \(d_k\to0\) gives \(g^T(z-x)=0\). Convexity
and the supporting inequality make \(F\) constant at \(\ell\) on
\([x,z]\), contrary to the strictly convex boundary of \(\{F\le\ell\}\).

Equations (16)--(17) and the same cutoff frame complete the proof. A
quadratic-plus-absolute-value objective is included, but those familiar
cases were already proved in earlier notes. The extension below
illustrates the additional normal irregularity allowed here.


\setcounter{section}{6}
\setcounter{theorem}{0}
\setcounter{lemma}{0}
\setcounter{corollary}{0}
\setcounter{proposition}{0}
\setcounter{remark}{0}
\setcounter{equation}{0}
\section{An explicit all-dimensional example outside
RS}\label{U:an-explicit-all-dimensional-example-outside-rs}

Let \(a_j=2^{-j}\) for \(j\ge2\), and let \[
 \eta(t)=\begin{cases}(1-t^2)^2,&|t|\le1,\\0,&|t|>1.\end{cases}
\] This nonnegative bump is \(C^1\), has \(\eta(0)=1\), and integral
\(16/15\). Its supports at the following scales are pairwise disjoint.
Define \[
 k(v)=\sum_{j=2}^{\infty}a_j^{-3}
 \eta\!\left(\frac{v-a_j}{a_j^5}\right),\qquad v>0,
\] \[
 r(v)=\begin{cases}
 \displaystyle\int_0^v(v-t)k(t)\,dt,&v>0,\\
 0,&v\le0.
 \end{cases}\tag{U.18}
\] Each bump has mass \((16/15)a_j^2\). The total mass is finite. The
sum is locally finite away from zero, so \(r\) is convex, \(C^1\)
globally, and \(C^3\) on \((0,\infty)\) and \((-\infty,0)\). In
particular it is locally \(C^{1,1}\) away from zero. Direct bounds give
\[
 0\le r'(v)\le6v^2,\qquad 0\le r(v)\le6v^3
\] near zero. However, \[
 r''(a_j)=a_j^{-3}\to\infty.\tag{U.19}
\] (The second derivative at zero itself exists and equals zero, but is
not continuous there.)

For any \(n\ge2\), consider \[
 \boxed{
 F(u,v)=\frac12\|u\|^2+\frac12v^2+|v|+r(v),
 \qquad u\in\mathbb R^{n-1}.
 }\tag{U.20}
\] It is 1-strongly convex and coercive, with unique minimizer zero. Its
only nondifferentiable hyperplane is \(v=0\); its tangential gradient is
\(u\), which is globally Lipschitz. Theorem U.1 therefore proves global
convergence of every admissible full-BFGS run, for arbitrary coupled SPD
initialization.

Yet \textbf{RS fails in every neighborhood of zero}. Fix any
sufficiently small nonzero \(u_0\), and take \(x_j=(u_0,a_j)\). Their
objective values approach the positive level \(\|u_0\|^2/2\). At these
smooth points, \[
 x_j^T\nabla^2F(x_j)x_j
 =\|u_0\|^2+a_j^2+a_j^{-1}\longrightarrow\infty.\tag{U.21}
\] A uniform inward Taylor estimate on that fixed positive band would,
by taking its second-order limit at each \(x_j\), bound the left side of
(21). This is impossible. The recorded equivalence between inward radial
smoothness and RS therefore rules out RS on that band. Choosing \(u_0\)
arbitrarily small gives failure in every neighborhood of the optimizer.

This is a \textbf{function-class} gain beyond RS, but not a claim that
the function lacks every prior transported certificate. It also expands
the allowed normal curvature rather than settling the original arbitrary
curved finite-\(C^2\)-max conjecture. Where each piece has a \(C^{1,1}\)
extension through the flat interface, the older
finite-polyhedral-partition theorem already gave local convergence.

\subsection{7.1 The infinite series is evaluated exactly in the
verifier}\label{U:the-infinite-series-is-evaluated-exactly-in-the-verifier}

There is no need to truncate (18). All bumps below a query can be summed
using geometric series, with at most one partially included bump. Define
\[
 P(t)=t-\frac23t^3+\frac15t^5,\qquad
 Q(t)=\frac12t^2-\frac12t^4+\frac16t^6.
\] For a query inside the bump centered at \(a=a_j\), with
\(t=(v-a)/a^5\), the partial mass is \(a^2[P(t)+8/15]\), and its first
moment is \[
 a^3[P(t)+8/15]+a^7[Q(t)-1/6].
\] All smaller bumps have total mass \((16/45)a^2\) and first moment
\((16/105)a^3\). This gives exact rational function and derivative
values for rational queries. At a peak, \[
 r(a)=\frac{64}{315}a^3+\frac16a^7,\qquad
 r'(a)=\frac89a^2,\qquad r''(a)=a^{-3}.
\] These are identities for the complete infinite series, not for a truncation. The exact rational-query formulas and their derivative values are proved for this oracle. For the objective (20), the convergence theorem applies in every stated dimension with an arbitrary coupled SPD initial inverse matrix; it does not replace full BFGS by a reduced two-dimensional or block-diagonal algorithm.

\subsection{7.2 The objective need not split as a smooth function plus
one scalar
penalty}\label{U:the-objective-need-not-split-as-a-smooth-function-plus-one-scalar-penalty}

For a fixed vector \(c\in\mathbb R^{n-1}\) with \(\|c\|<1\), another
example is \[
 F(u,v)=\frac12(\|u\|^2+v^2)+|(1+c^Tu)v|+r(v).
\] The two quadratic functions before adding \(r\) have Hessians \[
 \begin{pmatrix}I&c\\c^T&1\end{pmatrix},\qquad
 \begin{pmatrix}I&-c\\-c^T&1\end{pmatrix},
\] both bounded below by \((1-\|c\|)I\). Their maximum plus convex
\(r(v)\) is strongly convex. Since the absolute-value term and \(r\) are
nonnegative, zero is the unique minimizer.

On a neighborhood with \(|c^Tu|<1/2\), the only exceptional hyperplane
is \(v=0\), and the tangential gradient is \[
 A(u,v)=u+c|v|,
\] which is Lipschitz across that hyperplane. Theorem U.1 gives an objective-dependent local full-BFGS convergence basin for every $\|c\|<1$, including $c=0$.

If $c\ne0$, the normal-gradient jump magnitude $2(1+c^Tu)$ varies with $u$ on this neighborhood. More precisely, on no open neighborhood of zero can this objective equal $S(u,v)+h(v)$ with $S$ differentiable and $h$ an arbitrary scalar function. One way to see the obstruction is to subtract values at two nearby tangential points: the scalar term cancels, while the varying jump leaves a nonzero multiple of $|v|$, contradicting differentiability at $v=0$.

The nonzero condition is needed only for this nonsplitting conclusion. When $c=0$, the objective has the explicit global decomposition
\[
 S(u,v)=\tfrac12(\|u\|^2+v^2),\qquad h(v)=|v|+r(v),
\]
with $S$ smooth. Thus the subsection's existential claim is realized by nonzero $c$, not by every allowed parameter. The spike argument still makes (RS) fail in every neighborhood of zero for all $\|c\|<1$, including the zero case.


\endgroup
